\documentclass{article}
\usepackage[utf8]{inputenc}
\usepackage{amsmath}
\usepackage{amsfonts}
\usepackage{amssymb}

\title{Midterm Presentation Q\&A: MuscleWalker Project}
\author{Keval Shah}
\date{\today}

\begin{document}

\maketitle

\section{Project Overview}

\subsection{What is your project about?}

Our project aims to replace the traditional motor actuators in the Walker agent from the dm-control benchmark for reinforcement learning with artificial muscle tendons. The goal is to develop a bipedal walking agent that uses muscle-like actuation instead of conventional motors, and train effective policies for forward locomotion using either sampling-based Model Predictive Control (MPC) or learned reinforcement learning policies.

The project involves:
\begin{itemize}
    \item Creating a MuJoCo-based bipedal walker model with muscle tendon actuators
    \item Implementing control policies for forward walking locomotion
    \item Comparing performance between MPC and RL-based approaches
    \item Analyzing the biomechanical advantages of muscle-based actuation
\end{itemize}

\subsection{Why is this project important?}

This project is important for several reasons:

\textbf{Biomechanical Realism:} Artificial muscles provide more biologically realistic actuation compared to traditional motors, offering:
\begin{itemize}
    \item Non-linear force-length and force-velocity relationships
    \item Energy-efficient actuation through elastic energy storage
    \item More natural movement patterns that mimic biological locomotion
\end{itemize}

\textbf{Robotics Applications:} Muscle-based actuation has potential applications in:
\begin{itemize}
    \item Soft robotics and bio-inspired robot design
    \item Prosthetics and exoskeletons
    \item Humanoid robotics with more natural movement
\end{itemize}

\textbf{Control Theory:} Understanding how to control muscle-actuated systems advances:
\begin{itemize}
    \item Control algorithms for complex, non-linear actuators
    \item Learning-based control methods for biomechanical systems
    \item Hybrid control approaches combining MPC and RL
\end{itemize}

\subsection{How are you approaching this project?}

Our approach follows a systematic methodology:

\textbf{1. Model Development:}
\begin{itemize}
    \item Start with the dm-control Walker benchmark model
    \item Replace motor actuators with muscle tendon models in MuJoCo
    \item Implement proper muscle dynamics (Hill-type muscle model)
    \item Validate the biomechanical model against known muscle properties
\end{itemize}

\textbf{2. Control Policy Development:}
\begin{itemize}
    \item Implement sampling-based MPC for real-time control
    \item Develop RL-based policies using appropriate algorithms (PPO, SAC, etc.)
    \item Compare performance metrics: stability, energy efficiency, walking speed
\end{itemize}

\textbf{3. Evaluation Framework:}
\begin{itemize}
    \item Establish baseline performance with original motor actuators
    \item Compare muscle vs. motor performance across multiple metrics
    \item Analyze energy consumption and walking efficiency
\end{itemize}

\subsection{What have you accomplished so far?}

Based on the current repository state, we have accomplished:

\textbf{Infrastructure Setup:}
\begin{itemize}
    \item Created a complete MuJoCo-based bipedal walker simulation environment
    \item Implemented proper project structure with UV package management
    \item Set up visualization capabilities with MuJoCo's built-in viewer
    \item Created a functional walker model with 6 DOF (torso + 2 legs)
\end{itemize}

\textbf{Model Implementation:}
\begin{itemize}
    \item Developed a planar walker model with torso, thighs, shins, and feet
    \item Implemented 6 motor actuators (hip, knee, ankle for each leg)
    \item Added proper joint constraints and physical properties
    \item Created visual materials and rendering setup
\end{itemize}

\textbf{Basic Control:}
\begin{itemize}
    \item Implemented simple alternating leg control logic
    \item Created phase-based walking pattern (right leg forward, then left leg forward)
    \item Added position tracking and movement monitoring
    \item Established real-time simulation with visual feedback
\end{itemize}

\textbf{Technical Foundation:}
\begin{itemize}
    \item Python 3.13+ environment with MuJoCo 3.0+ integration
    \item Proper dependency management with UV
    \item Cross-platform compatibility (macOS, Linux, Windows)
    \item Interactive 3D visualization with camera controls
\end{itemize}

\subsection{What do you plan to do next?}

Our immediate next steps include:

\textbf{Week 1: Muscle Implementation}
\begin{itemize}
    \item Replace motor actuators with Hill-type muscle models
    \item Implement basic muscle-tendon dynamics with force-length relationships
    \item Add muscle activation dynamics
    \item Validate muscle model functionality
\end{itemize}

\textbf{Week 2: Control Development}
\begin{itemize}
    \item Implement sampling-based MPC controller
    \item Develop RL environment with reward functions
    \item Train initial RL policies using PPO/SAC
    \item Optimize control parameters for muscle-actuated walking
\end{itemize}

\textbf{Week 3: Analysis and Optimization}
\begin{itemize}
    \item Compare muscle vs. motor performance metrics
    \item Analyze energy efficiency and walking stability
    \item Fine-tune control policies for better performance
    \item Document key findings and challenges
\end{itemize}

\textbf{Week 4: Documentation and Finalization}
\begin{itemize}
    \item Create technical documentation
    \item Prepare presentation materials
    \item Write project report
    \item Finalize code and prepare for submission
\end{itemize}

\textbf{Technical Milestones:}
\begin{itemize}
    \item Successful muscle-actuated walking without falling
    \item Stable forward locomotion for extended periods
    \item Demonstrated energy efficiency improvements over motors
    \item Robust control policies that generalize across different conditions
\end{itemize}

\end{document}
